\documentclass[12pt]{article}
\usepackage[utf8]{inputenc}
\usepackage[T1]{fontenc}
\usepackage[margin=1in]{geometry}
\usepackage{setspace}
\usepackage{xcolor}
\usepackage[colorlinks=true, urlcolor=blue!60!black, linkcolor=blue!60!black]{hyperref}
\setlength{\parindent}{0em}
\setlength{\parskip}{0.45em}
\renewcommand{\baselinestretch}{1.0}
\begin{document}
\begin{center}
{\Large \textit{Myers v. United States}, 272 U.S. 52 (1926)\par}
\vspace{0.3em}
{\normalsize Case Brief --- One-Page Edition\par}
\end{center}

\textbf{Memory Jogger.} The postmaster removal case and the foundation of presidential removal power (\href{https://supreme.justia.com/cases/federal/us/272/52/}{272 U.S. 52}).

\textbf{Facts.} President Wilson removed Portland postmaster Frank Myers before his four-year term ended without the Senate consent an 1876 statute required. Myers sued for his lost salary.

\textbf{Procedural History.} The Court of Claims ruled for the government; the Supreme Court affirmed on constitutional grounds after reargument.

\textbf{Votes.} 6--3. Taft wrote for the Court; Holmes, McReynolds, and Brandeis dissented separately.

\textbf{Holding.} Article II gives the President exclusive power to remove executive officers he appoints with Senate consent; Congress may not require Senate approval of removals.

\textbf{Reasoning.} The Vesting and Take Care Clauses imply control of subordinates; the Senate's consent role is a narrow exception; and the Decision of 1789 confirmed a constitutional removal power. Holmes and McReynolds said Congress may condition tenure in offices it creates. Brandeis replied that separation of powers prevents arbitrary power rather than promoting efficiency. \textit{Humphrey's Executor} narrowed \textit{Myers}; \textit{Trump v.\ Slaughter} (2026) revived it.

\textbf{Examples.} \emph{Same-Side:} for-cause protection for executive commissioners. \emph{Opposite-Side:} a department-appointed inferior officer (\textit{Perkins}). \emph{Fence-Sitter:} single-layer protection for an agency's inspector.

\textbf{Critique.} Historians say Taft overstated the 1789 consensus.

\textbf{Key Quotation.} ``The doctrine of the separation of powers was adopted by the convention of 1787, not to promote efficiency but to preclude the exercise of arbitrary power. The purpose was not to avoid friction, but, by means of the inevitable friction incident to the distribution of the governmental powers among three departments, to save the people from autocracy.'' --- Justice Brandeis, dissenting, 272 U.S.\ at 293.

\end{document}
